% Part: computability
% Chapter: recursive-functions
% Section: pr-functions

\documentclass[../../../include/open-logic-section]{subfiles}

\begin{document}

\olfileid[hi]{cmp}{rec}{prf}
\olsection{आदिम पुनरावर्ती फलन}

आइए फिर याद करें कि आदिम पुनरावर्तन और संयोजन का उपयोग करके हम विद्यमान
फलनों से नए फलन कैसे परिभाषित कर सकते हैं।

\begin{defn}
  \ollabel{defn:primitive-recursion} मान लें कि $f$ कोई $k$-स्थानी
  फलन ($k\ge 1$) है और $g$ कोई $(k+2)$-स्थानी फलन है। \emph{$f$ और
  $g$ से आदिम पुनरावर्तन द्वारा} परिभाषित फलन वह $(k+1)$-स्थानी
  फलन~$h$ है जिसे निम्न समीकरण परिभाषित करते हैं:
  \begin{align*}
    h(x_0,\dots,x_{k-1},0) & =  f(x_0,\dots,x_{k-1}) \\
    h(x_0,\dots,x_{k-1},y+1) & = g(x_0,\dots,x_{k-1}, y, h(x_0,\dots,x_{k-1}, y))
  \end{align*}
\end{defn}

\begin{defn}
  \ollabel{defn:composition} मान लें कि $f$ कोई $k$-स्थानी फलन है और
  $g_0$, \dots, $g_{k-1}$ ऐसे $k$ फलन हैं जो सभी $n$-स्थानी हैं।
  \emph{$f$ और $g_0$, \dots,~$g_{k-1}$ से संयोजन द्वारा} परिभाषित
  फलन वह $n$-स्थानी फलन~$h$ है जिसे निम्न समीकरण परिभाषित करते हैं:
  \[
  h(x_0, \dots, x_{n-1}) =
  f(g_0(x_0, \dots, x_{n-1}), \dots, g_{k-1}(x_0, \dots, x_{n-1})).
  \]
\end{defn}

$\Succ$ और प्रक्षेपण फलनों
\[
\Proj{n}{i}(x_0,\dots,x_{n-1}) = x_i,
\]
के अतिरिक्त, प्रत्येक प्राकृतिक संख्या $n$ और $i < n$ के लिए हम
फलन~$\Zero(x) = 0$ को आदिम पुनरावर्ती फलनों में सम्मिलित करेंगे।

\begin{defn}
  आदिम पुनरावर्ती फलनों का समुच्चय, $\Nat^n$ से $\Nat$ में फलनों का वह
  समुच्चय है जिसे निम्न उपवाक्यों से आगमनात्मक रूप से परिभाषित किया गया है:
  \begin{enumerate}
  \item $\Zero$ आदिम पुनरावर्ती है।
  \item $\Succ$ आदिम पुनरावर्ती है।
  \item प्रत्येक प्रक्षेपण फलन $\Proj{n}{i}$ आदिम पुनरावर्ती है।
  \item यदि $f$ कोई $k$-स्थानी आदिम पुनरावर्ती फलन है और $g_0$,
    \dots,~$g_{k-1}$, $n$-स्थानी आदिम पुनरावर्ती फलन हैं, तो $f$ का
    $g_0$, \dots,~$g_{k-1}$ के साथ संयोजन आदिम पुनरावर्ती है।
  \item यदि $f$ कोई $k$-स्थानी आदिम पुनरावर्ती फलन है और $g$ कोई
    $k+2$-स्थानी आदिम पुनरावर्ती फलन है, तो $f$ और $g$ से आदिम
    पुनरावर्तन द्वारा परिभाषित फलन आदिम पुनरावर्ती है।
\end{enumerate}
\end{defn}

\begin{explain}
अधिक संक्षेप में, आदिम पुनरावर्ती फलनों का समुच्चय वह लघुतम समुच्चय है
जिसमें $\Zero$, $\Succ$ और प्रक्षेपण फलन~$\Proj{n}{j}$ आते हैं और जो
संयोजन तथा आदिम पुनरावर्तन के अधीन संवृत है।

आदिम पुनरावर्ती फलनों के समुच्चय का वर्णन करने का एक अन्य ढंग उसे
``चरणों'' के पदों में परिभाषित करना है। $S_0$ से आरंभिक फलनों का समुच्चय
निरूपित करें: $\Zero$, $\Succ$ और प्रक्षेपण। ये चरण~$0$ के आदिम
पुनरावर्ती फलन हैं। किसी चरण $S_i$ के परिभाषित हो जाने पर $S_{i+1}$ को
उन सभी फलनों का समुच्चय मानें जो~$S_i$ में पहले से उपस्थित फलनों पर
संयोजन अथवा आदिम पुनरावर्तन का एक अनुप्रयोग करके मिलते हैं। तब
\[
S = \bigcup_{i \in \Nat} S_i
\]
सभी आदिम पुनरावर्ती फलनों का समुच्चय है।
\end{explain}

आइए सत्यापित करें कि $\Add$ एक आदिम पुनरावर्ती फलन है।

\begin{prop}
  जोड़ फलन $\Add(x,y) = x+y$ आदिम पुनरावर्ती है।
\end{prop}

\begin{proof}
हमारे पास $\Add$ की आदिम पुनरावर्ती परिभाषा दो फलनों $f$ और~$g$ के
पदों में पहले से है, जो \olref{defn:primitive-recursion} के रूप से मेल
खाती है:
\begin{align*}
  \Add(x_0, 0) & = f(x_0) = x_0\\
  \Add(x_0, y+1) & = g(x_0, y, \Add(x_0, y)) =
  \Succ(\Add(x_0, y))
\end{align*}
$\Add$ का आदिम पुनरावर्ती होना इस पर निर्भर है कि $f$ और $g$ भी आदिम
पुनरावर्ती हों। $f(x_0) = x_0 = \Proj{1}{0}(x_0)$ है और प्रक्षेपण फलन आदिम
पुनरावर्ती माने जाते हैं, इसलिए $f$ आदिम पुनरावर्ती है। फलन $g$ वह
त्रि-स्थानी फलन $g(x_0, y, z)$ है जिसे इस प्रकार परिभाषित किया गया है:
\[
g(x_0, y, z) = \Succ(z).
\]
इससे अभी यह नहीं मिलता कि $g$ आदिम पुनरावर्ती है, क्योंकि $g$ और
$\Succ$ ठीक एक ही फलन नहीं हैं: $\Succ$ एक-स्थानी है, जबकि $g$ को
त्रि-स्थानी होना है। किंतु हम संयोजन द्वारा $g$ को ``औपचारिक रूप से''
इस प्रकार परिभाषित कर सकते हैं:
\[
g(x_0, y, z) = \Succ(\Proj{3}{2}(x_0, y, z))
\]
चूँकि $\Succ$ और $\Proj{3}{2}$ आदिम पुनरावर्ती फलन माने जाते हैं, इसलिए
$g$ भी ऐसा ही है, क्योंकि उसे आदिम पुनरावर्ती फलनों से संयोजन द्वारा
परिभाषित किया जा सकता है।
\end{proof}

\begin{prop}
  \ollabel{prop:mult-pr}
  गुणा फलन $\Mult(x,y) = x \cdot y$ आदिम पुनरावर्ती है।
\end{prop}

\begin{proof}
  अभ्यास।
\end{proof}

\begin{prob}
  यह दिखाकर \olref[cmp][rec][prf]{prop:mult-pr} सिद्ध करें कि $\Mult$
  की आदिम पुनरावर्ती परिभाषा को
  \olref[cmp][rec][prf]{defn:primitive-recursion} द्वारा अपेक्षित रूप
  में रखा जा सकता है और संगत फलन $f$ तथा~$g$ आदिम पुनरावर्ती हैं।
\end{prob}

\begin{ex}
आदिम पुनरावर्ती परिभाषा का हमारा सबसे पहला उदाहरण यह है:
\begin{align*}
h(0) & =  1 \\
h(y+1) & =  2 \cdot h(y).
\end{align*}
यह फलन \olref{defn:primitive-recursion} द्वारा अपेक्षित रूप में नहीं आ
सकता, क्योंकि $k=0$ है। परिभाषा में अचर $1$ और $2$ भी आते हैं। पहली
  समस्या से निपटने के लिए एक निष्क्रिय तर्क जोड़कर फलन~$h'$ परिभाषित करें:
\begin{align*}
h'(x_0, 0) & =  f(x_0) = 1 \\
h'(x_0, y+1) & =  g(x_0, y, h'(x_0, y)) = 2 \cdot h'(x_0, y).
\end{align*}
फलन $f(x_0) = 1$ को $\Succ$ और $\Zero$ से संयोजन द्वारा परिभाषित किया
जा सकता है: $f(x_0) = \Succ(\Zero(x_0))$। फलन $g$ को
$g'(z) = 2 \cdot z$ और प्रक्षेपणों से संयोजन द्वारा परिभाषित किया जा
सकता है:
\begin{align*}
g(x_0, y, z) & = g'(\Proj{3}{2}(x_0, y, z))
\intertext{और फिर $g'$ को संयोजन द्वारा इस प्रकार परिभाषित किया जा सकता है:}
g'(z) & = \Mult(g''(z), \Proj{1}{0}(z))
\intertext{और}
g''(z) & = \Succ(f(z)),
\end{align*}
जहाँ $f$ ऊपर वाला ही है: $f(z) = \Succ(\Zero(z))$। अब हमारे पास~$h'$
है, इसलिए फिर संयोजन का उपयोग करके $h(y) =
h'(\Proj{1}{0}(y),\Proj{1}{0}(y))$ रखें। इससे पता चलता है कि $h$ को मूल
फलनों से संयोजनों और आदिम पुनरावर्तनों के एक अनुक्रम द्वारा परिभाषित किया
जा सकता है; अतः $h$ आदिम पुनरावर्ती है।
\end{ex}

\end{document}
